% GENERATED FILE. Edit canonical lessons, then run npm run generate:books.
% canonical-source-hash: fnv1a64:1af52d0a4b5e25f2
% canonical-lessons: ES-C762-crucigrama, ES-C762-revista, ES-C762-lector, ES-C762-portada, ES-C762-hojear

\chapter{Leer en casa}
\label{ch:es-leer-en-casa}
\hlchaptermodality{762}

\begin{chapteropening}
\textbf{By the end of this chapter:} \emph{I can talk about reading at home: magazines, covers, crosswords and flicking through pages.}

Complete the last lesson of chapter 762: I can talk about reading at home: magazines, covers, crosswords and flicking through pages.
\end{chapteropening}

\section[el crucigrama]{el crucigrama — the crossword}

\label{lesson:ES-C762-crucigrama}



\begin{quote}
Before the new one: say the Spanish for to yawn, then the Spanish for to record.
\end{quote}

\subsection*{You'll want to know: el crucigrama}
\textbf{el crucigrama} — \textquotedblleft{}the crossword\textquotedblright{}.

\begin{quote}
Hago el crucigrama del periódico.
\end{quote}

\begin{quote}
\emph{I do the newspaper crossword.}
\end{quote}

Masculine, though it ends in \emph{-a}.

\begin{cousinweb}
\emph{crucigrama} is built on \emph{cruz} \textquotedblleft{}cross\textquotedblright{}, from Latin \emph{crux}, and Greek \emph{grámma} \textquotedblleft{}letter\textquotedblright{}, a twentieth-century coinage for crossed words.
\end{cousinweb}

\begin{grammarlens}[title={What you've built}]
One more word for telling, explaining and describing at B1.
\end{grammarlens}

\subsection*{Your turn}
\begin{itemize}
\raggedright
  \item \emph{Say it:} \emph{el crucigrama}
  \item \emph{Say it:} \emph{el crucigrama}, once more
  \item \emph{Say it:} a sentence of your own with \emph{el crucigrama}
  \item \emph{Recall:} say the Spanish for to yawn, then the Spanish for to record, then say \emph{el crucigrama} again
\end{itemize}

\begin{tcolorbox}[breakable,colback=teal!4,colframe=teal!35!black,title={Before you move on}]
What does el crucigrama mean? (\textquotedblleft{}The crossword\textquotedblright{}.) Say it once more.
\end{tcolorbox}
\section[la revista]{la revista — the magazine}

\label{lesson:ES-C762-revista}



\begin{quote}
Before the new one: say the Spanish for to record, then the Spanish for the crossword.
\end{quote}

\subsection*{You'll want to know: la revista}
\textbf{la revista} — \textquotedblleft{}the magazine\textquotedblright{}.

\begin{quote}
Compro una revista de cocina.
\end{quote}

\begin{quote}
\emph{I buy a cooking magazine.}
\end{quote}

\begin{cousinweb}
\emph{revista} is built on \emph{re-} and \emph{vista} \textquotedblleft{}view\textquotedblright{}, from Latin \emph{vidēre} \textquotedblleft{}to see\textquotedblright{}; it was probably modelled on English \emph{review} and French \emph{revue}.
\end{cousinweb}

\begin{grammarlens}[title={What you've built}]
One more word for telling, explaining and describing at B1.
\end{grammarlens}

\subsection*{Your turn}
\begin{itemize}
\raggedright
  \item \emph{Say it:} \emph{la revista}
  \item \emph{Say it:} \emph{la revista}, once more
  \item \emph{Say it:} a sentence of your own with \emph{la revista}
  \item \emph{Recall:} say the Spanish for to record, then the Spanish for the crossword, then say \emph{la revista} again
\end{itemize}

\begin{tcolorbox}[breakable,colback=teal!4,colframe=teal!35!black,title={Before you move on}]
What does la revista mean? (\textquotedblleft{}The magazine\textquotedblright{}.) Say it once more.
\end{tcolorbox}
\section[el lector]{el lector — the reader}

\label{lesson:ES-C762-lector}



\begin{quote}
Before the new one: say the Spanish for the crossword, then the Spanish for the magazine.
\end{quote}

\subsection*{You'll want to know: el lector}
\textbf{el lector} — \textquotedblleft{}the reader\textquotedblright{}.

\begin{quote}
Es un gran lector de novelas.
\end{quote}

\begin{quote}
\emph{He is a great reader of novels.}
\end{quote}

\emph{La lectora} for a woman. \emph{Un lector de libros electrónicos} — an e-reader.

\begin{cousinweb}
\emph{lector} is a learned borrowing from Latin \emph{lēctor} \textquotedblleft{}reader\textquotedblright{} (EtymDB), from \emph{legere} \textquotedblleft{}to read\textquotedblright{}.
\end{cousinweb}

\begin{grammarlens}[title={What you've built}]
One more word for telling, explaining and describing at B1.
\end{grammarlens}

\subsection*{Your turn}
\begin{itemize}
\raggedright
  \item \emph{Say it:} \emph{el lector}
  \item \emph{Say it:} \emph{el lector}, once more
  \item \emph{Say it:} a sentence of your own with \emph{el lector}
  \item \emph{Recall:} say the Spanish for the crossword, then the Spanish for the magazine, then say \emph{el lector} again
\end{itemize}

\begin{tcolorbox}[breakable,colback=teal!4,colframe=teal!35!black,title={Before you move on}]
What does el lector mean? (\textquotedblleft{}The reader\textquotedblright{}.) Say it once more.
\end{tcolorbox}
\section[la portada]{la portada — the cover; the front page}

\label{lesson:ES-C762-portada}



\begin{quote}
Before the new one: say the Spanish for the magazine, then the Spanish for the reader.
\end{quote}

\subsection*{You'll want to know: la portada}
\textbf{la portada} — \textquotedblleft{}the cover; the front page\textquotedblright{}.

\begin{quote}
Me gusta la portada de esta revista.
\end{quote}

\begin{quote}
\emph{I like the cover of this magazine.}
\end{quote}

\emph{En portada} — on the front page.

\begin{cousinweb}
\emph{portada} is usually traced to Latin \emph{porta} \textquotedblleft{}door, gate\textquotedblright{}, the source of \emph{puerta}: it named the front of a building, then the cover of a book.
\end{cousinweb}

\begin{grammarlens}[title={What you've built}]
One more word for telling, explaining and describing at B1.
\end{grammarlens}

\subsection*{Your turn}
\begin{itemize}
\raggedright
  \item \emph{Say it:} \emph{la portada}
  \item \emph{Say it:} \emph{la portada}, once more
  \item \emph{Say it:} a sentence of your own with \emph{la portada}
  \item \emph{Recall:} say the Spanish for the magazine, then the Spanish for the reader, then say \emph{la portada} again
\end{itemize}

\begin{tcolorbox}[breakable,colback=teal!4,colframe=teal!35!black,title={Before you move on}]
What does la portada mean? (\textquotedblleft{}The cover; the front page\textquotedblright{}.) Say it once more.
\end{tcolorbox}
\section[hojear]{hojear — to leaf through, to flick through}

\label{lesson:ES-C762-hojear}



\begin{quote}
Before the new one: say the Spanish for the reader, then the Spanish for the cover.
\end{quote}

\subsection*{You'll want to know: hojear}
\textbf{hojear} — \textquotedblleft{}to leaf through, to flick through\textquotedblright{}.

\begin{quote}
Hojeo una revista en el tren.
\end{quote}

\begin{quote}
\emph{I flick through a magazine on the train.}
\end{quote}

Not to be confused with \emph{ojear}, to glance at.

\begin{cousinweb}
\emph{hojear} is built on \emph{hoja} \textquotedblleft{}leaf, page\textquotedblright{} with the ending \emph{-ear}; \emph{hoja} is inherited from Latin \emph{folia} \textquotedblleft{}leaves\textquotedblright{} (EtymDB), with \emph{f-} turned to \emph{h-}.
\end{cousinweb}

\begin{grammarlens}[title={What you've built}]
One more word for telling, explaining and describing at B1.
\end{grammarlens}

\subsection*{Your turn}
\begin{itemize}
\raggedright
  \item \emph{Say it:} \emph{hojear}
  \item \emph{Say it:} \emph{hojear}, once more
  \item \emph{Say it:} a sentence of your own with \emph{hojear}
  \item \emph{Recall:} say the Spanish for the reader, then the Spanish for the cover, then say \emph{hojear} again
\end{itemize}

\begin{tcolorbox}[breakable,colback=teal!4,colframe=teal!35!black,title={Before you move on}]
What does hojear mean? (\textquotedblleft{}To leaf through, to flick through\textquotedblright{}.) Say it once more.
\end{tcolorbox}
